\part{Unchecked targets and remaining gates}
\label{part:open}
\section{What has not been proved}
There is no unchecked finding offered as a theorem in this part. The
following are the unresolved targets of the task. Their state is a
current gate, not a proved negative answer. A precise counterexample to
one proposed transfer is recorded separately from a proof that no other
transfer can exist.

\begin{description}
\item[O-PH2-SH: direct bounds for irreducible-inner plethysm.]
Both clauses of ordinary plethysm PH3 are now refuted: P-F19 gives
an exponential canonical parameter degree, and P-F30--P-F31 prove
actual simple poles of exponential order in the same cubic family.
The complete support classification closes the singular-boundary
and possible-cancellation gates of P-F23--P-F29. The first PH3
clause therefore cannot serve as the premise for a proof of PH2,
although the conditional implication P-F11--P-F12 remains valid.

The unresolved question is a direct proof or refutation of the
narrower polynomial positivity-index and saturation-index bounds,
with the source's exact parameter restrictions and constituent
polynomial definitions. A large quasipolynomial period does not by
itself force either index to be large. P-F32 supplies a precise
exterior-power complement identity and an infinite family with
first positive stretch $\Theta(\sqrt n)$; it gives only a polynomial
index lower bound. It does not settle SH(a), its separate density
clause SH(b), or PH2. P-F14--P-F15 already refute the broader
rank-only subgroup-restriction versions, a different formulation.

P-F34 makes a direct sufficient gate precise. Let $B_n$ be the
average constituent polynomial, determined by the principal part
of $H_n$ at one, and let $U_n$ be its coefficient $\ell^1$ norm.
For a specified primitive root, the proved Fourier amplitude satisfies
$\log_2|A_n|=n^4/24+O(n^3)$. Positivity or saturation
after a shift forces $B_n$ at that shift (or one more) to dominate
this amplitude. The additional bound
$\limsup (\log_2 U_n)/n^4<1/24$ would therefore refute polynomial
PH2 and SH(a). This bound remains unproved; the Fourier inequality
is necessary, not an equivalence to the index conjectures.
P-F38 excludes deriving this bound from small matrix entries, bounded
coordinates and cubic-scale dimension alone, even after normalizing by
the row rank. The actual cubic weight system and Weyl alternation remain
necessary extra structure. P-F33 shows why normal rational Gorenstein geometry with fixed
negative $a$-invariant and fixed dimension cannot alone supply
such a bound: its genuine arbitrarily large periods coexist with
zero indices and an unbounded average constant term. P-F35 gives
the exact calibration $p(f_4)=s(f_4)=23$ and minimal period $120$.

\item[O-COUNT: uniform effective shifted counting.]
P-F2 constructs a single shifted polytope from explicit positive Hilbert
data. P-F12 controls their existence and size under a bounded
positive-form premise in expanded parameters. P-F31 shows that the
proposed numerical PH3 bound cannot supply that premise in general. P-F13 reconstructs
and verifies candidate data from a finite prefix under an explicit
denominator promise. A procedure constructing them with the original
partition-input complexity bounds remains to be given. A small object
whose existence is known separately for every instance is not automatically
a polynomial-time constructible object. The homogeneous version is already
refuted by P-K1; no revision should silently restore that false assertion.

\item[O-NSQ: nonstandard geometric realization.]
For the particular bases and products in GCT VII--VIII, one must construct
the category, identify the basis, control parity or shifts, establish
product and specialization compatibility, and supply size bounds.
P-F6 is a sufficient criterion after these data exist. The determinant
and theta-lift passages of the supplied note do not themselves identify
the relevant Hom spaces. P-F5 refutes the simplest attempted substitution
of ordinary interlacing for plethysm restriction.
P-F36 additionally proves that the actual cubic example has a
quadratic coefficient extension that survives every finite power
regrading. Its separate conjugate blocks cannot descend through
an ordinary rational projector, and its displayed raw coefficient
cannot be a finite graded Laurent character. A comparison must
retain that extension and its descent datum, or explicitly recover
the blocks from a conjugate sum. Arbitrary coefficient systems
and the source's relaxed minimal-polynomial conjectures remain
outside this obstruction.

\item[O-OBSTRUCTION: an actual family of separating equations.]
In the no-occurrence range, prove a strict inequality between the actual
coordinate-ring multiplicities, or construct an explicit polynomial
vanishing on the determinant closure and not on the padded permanent.
The ambient plethysm and Kronecker multiplicities do not determine the
ranks of the quotient maps. P-F8 and PL-06 isolate this rank/comparison
requirement. No such asymptotic obstruction family is constructed here.

\item[O-DIR-HN: arbitrary solution fields in Nullstellensatz.]
C-GRH-01 proves only the cyclotomic-field transfer. Koiran's algorithm uses
a uniform prime-ideal/root-modulo-prime estimate for arbitrary primitive
solution fields. Rojas's DZH or GIPIT conditions are alternative sufficient
analytic inputs. The present supplied Dirichlet theorem does not state
these inputs. No formal nonimplication between the analytic conjectures
is proved.

\item[O-PIT: the derandomization required by GCT V.]
The exact algorithm or hardness hypothesis of each theorem must still be
provided. The supplied counting algorithm yields C-PIT-01 in a bounded
finite-field circuit model. It does not supply the characteristic-zero
black-box, border or strengthened symbolic determinant PIT assertions
used elsewhere. GRH does not replace those separate hypotheses.

\item[O-INTEGRAL: uniform bounds for specialization.]
B7.1 proves good-prime survival for a fixed integral separator; B7.2 gives
an exact rank test in any fixed degree and integer minors excluding bad
primes. An efficient GCT transfer needs a uniform separator-degree bound,
succinct rank data and a usable prime bound. Neither Hilbert-scheme
representability nor the pointwise specialization theorem supplies these
quantitative statements.

\item[O-COMPLEXITY: Boolean and algebraic separation.]
The ETH theorem refutes the precise exponential hypotheses treated in
C-SAT-01. Its supplied circuit corollary is superlinear for
$\mathsf E^{\mathsf{NP}}$, not superpolynomial for $\mathsf{NP}$.
The current constant-free integer Boolean--arithmetic disjunction cited
in the complexity section does not choose a branch and does not identify
the usual complex-field $\mathsf{VP}$ with its constant-free integer
variant. The stated upper bounds and conditional bridges do not settle
$\mathsf P$ versus $\mathsf{NP}$ or the target GCT obstruction.
\end{description}

\section{Approaches tried and the next discriminating steps}
The qualitative-sign route was tested by actual determinant projectors
and a virtual-character counterexample (P-F3--P-F4). Its next step is an
identity with the intended multiplicity projector together with a size
bound. The standard-branching route was tested on $\operatorname{Sym}^2$
(P-F5); a replacement must use the actual Schur-functor homomorphism.

The purity route was tested at three increasingly specific levels:
semisimple categorical structure (P-F6), pure finite schemes encoding
\#3SAT (P-F7), and the actual canonical plethysm algebra carrying the
obstructing Hilbert function, a positive multiplicative basis and pure
intersection cohomology (PLE-F7--PLE-F8). The next step is an
equivariant graded comparison theorem, not another proof that an unrelated
cohomology group is pure. P-F10 separately tests the missing quantitative
control of a grading.

The Hilbert-series route made positive progress: P-F1 constructs the
optimal shifted model on the obstructing ray; PLE-F5 realizes its normal
invariant ring; PLE-F6 classifies all rescalings that repair homogeneous
Ehrhart realization; P-F2 treats every explicit positive rational form;
P-F9 and P-F11 supply coefficient-positivity bounds; P-F12 connects the
last bound to the numerical positive-form clause of revised PH3.
P-F13 gives finite reconstruction under a denominator promise;
P-F16--P-F18 give finite fixed-embedding and fixed-dimension bounds;
P-F19 refutes polynomial growth of the canonical parameter bound.
P-F30--P-F31 subsequently close the numerical positive-form question
negatively, using a complete support theorem and a fixed-component
calculation. The next step is a direct PH2 or saturation argument,
together with uniform construction of a suitable revised model.
The ordinary pole is proved by these additional arguments, not
inferred from the h.s.o.p. bound alone.
P-F20 confirms that distinction with universal denominators and
fixed-target covariants. Its multivariate counterexample defeats a
transfer based only on quasi-polynomiality from an affine slice to
simultaneous rays. P-F21--P-F22 then defeat generic normal,
Cohen--Macaulay, rational and Gorenstein attempts to repair the
parameter-degree inference. P-F23--P-F24 made the geometric gate precise: the known chain
component is smooth, and the other fixed components had to be
controlled. P-F26 identifies the entire smooth fixed locus with
the chain orbit; P-F27 excludes every graph-support boundary
competitor, including rank-deficient supports and lifts of larger
two-power order. P-F29 handles the complementary cube-and-squarefree
case. P-F30 excludes all remaining mixed supports, and P-F31 applies
the noncancellation criterion. No competing nonzero component remains. The binary determinant subfamily is separately treated
in P-F25: its realized holes are bounded by thirteen, but the
source's strict saturation definition also tests off-residue
constituent values and is not settled by that table.

The archive route distinguishes the three products on representations,
separates the two Burnside operations, checks the transcribed identities
on small groups, and treats the modular and quadratic-form caveats
(PL-01--PL-08). Its obstruction gate is the actual rank of the
Foulkes--Howe multiplicity-space map. A Schur-positive character of its
source and target does not compute that rank.

The norm and approximation routes were tested against scaling, exact
zero detection, adjacent integer multiplicities, and an explicit rank-two
family approaching rank one in normalized norm (B4.1--B5.3). They do give
valid contraction-model obstructions under the stated operator bounds.
A full algebraic-orbit obstruction requires an invariant and a certified
separating gap that survive the allowed degenerations.

The arithmetic and ETH routes track the exact analytic family, all input
and reduction sizes, the direction of reductions, arithmetic precision,
and the class in which a lower bound is obtained. Their next steps are
specified in O-DIR-HN, O-PIT and O-COMPLEXITY. These are not resolved by
relabeling the supplied subexponential bound as polynomial time.
